\documentclass[10pt]{article}
\usepackage[margin=0.66in]{geometry}
\usepackage[T1]{fontenc}
\usepackage[utf8]{inputenc}
\usepackage{xcolor}
\usepackage{booktabs}
\usepackage{array}
\usepackage{tabularx}
\usepackage[hidelinks]{hyperref}
\setlength{\parindent}{0pt}
\setlength{\parskip}{0.26em}
\setlength{\emergencystretch}{2em}

\newcommand{\questionsep}{\par\vspace{0.75em}}

\begin{document}

\begin{center}
{\Large Response to Reviewer vCPF}\\[2pt]
Submission 13175
\end{center}

\textbf{Opening response.}
We sincerely thank the reviewer. The three requests sharpen the revision: audit
Visium HD construction, demonstrate breadth, and biologically calibrate low
correlations. This resource has accompanied me since my first PhD year; near
graduation, I hope its release helps newcomers enter spatial transcriptomics.

Our 21-sample in-house DBiC-seq collection adds 54,304 measured/53,989 post-QC
cells with paired morphology, RNA, and context. The working manifest contains 99
unique samples (100 records) and 28,315,247 targets; a broader continuing
DBiC-seq pool of approximately 200,000 paired cells is not counted here.

\questionsep
\textbf{Q1.}\quad
\textit{For the Visium HD data, how reliable is the bin-to-cell aggregation?
How are boundary bins assigned, and do segmentation errors noticeably affect
expression profiles?}

\textbf{A1 (Construction).}
\begin{itemize}
\item We will distinguish native Xenium cells from \emph{derived cell-aligned}
Visium HD targets.
\item Official transforms register native 2-\(\mu\)m bins to H\&E. A
closed bin square is assigned when it intersects a same-frame CellViT polygon;
if several polygons claim it, the nearest full-resolution centroid wins, with
raw-cell index breaking exact ties.
\item Unsupported cells are removed. We will move this rule and a complete
record example from Appendix S2 into the main text.
\end{itemize}

\textbf{A2 (Sensitivity audit).}\enspace We audited 3,000 raw polygons per
dataset (9,000 total) under 1-\(\mu\)m
registration shifts and mask erosion/dilation:

\begin{center}
\begingroup
\scriptsize
\setlength{\tabcolsep}{2pt}
\begin{tabularx}{\linewidth}{@{}Xrrrr@{}}
\toprule
Visium HD sample & Raw polygons with canonical bins & Shift-bin Jaccard &
Shift-expression cosine & Shift median \(|\Delta\mathrm{UMI}|\) \\
\midrule
Human lung cancer & 49.4\% & 0.727--0.733 & 0.954--0.957 & 11.1\%--11.3\% \\
Mouse brain & 97.5\% & 0.806--0.816 & 0.936--0.939 & 6.0\%--6.1\% \\
Human pancreas & 50.0\% & 0.706--0.714 & 0.994 & 13.0\%--13.7\% \\
\bottomrule
\end{tabularx}
\endgroup
\end{center}

\textbf{A3 (Interpretation).}
\begin{itemize}
\item Percentages use all detected polygons; unsupported polygons are excluded.
\item Shifts preserve expression direction but alter membership/UMIs;
erosion/dilation gives Jaccard 0.462--0.720, cosine 0.871--0.993, and
\(|\Delta\mathrm{UMI}|\) 32.8\%--66.6\%.
\item Uncertainty is material. We will release settings/QC, avoid ground-truth
language, and report Xenium separately.
\end{itemize}

\questionsep
\textbf{Q2.}\quad
\textit{The main experiments cover only a few representative cases. Can the
authors summarize more organs or platforms?}

\textbf{A1 (Release-wide breadth).}\enspace
Yes. Table~1 reports results for all 25 release categories.

\newpage
\begin{center}
\begingroup
\footnotesize
\setlength{\tabcolsep}{4.2pt}
\renewcommand{\arraystretch}{1.04}
\textit{\textbf{Table 1. Per-organ benchmark results.}
Results use the top-50 training-selected HVGs, four contiguous spatial
holdouts, and a 5\% train--test buffer.}\\[5pt]
\begin{tabular*}{\linewidth}{@{\extracolsep{\fill}}lrrrrrr@{}}
\toprule
& \multicolumn{3}{c}{Gene Pearson} & & & \\
\cmidrule(lr){2-4}
\shortstack[l]{Release category\\(evaluated species)} &
Image &
\shortstack{Coordinate\\only} &
\shortstack{Spatial\\KNN} &
\shortstack{Gene\\Spearman} &
\shortstack{Cell\\Pearson} &
F1 \\
\midrule
\textbf{Bone (mouse)} & 0.030 & \textbf{0.124} & 0.055 & 0.037 & 0.180 & 0.240 \\
Brain (human + mouse) & \textbf{0.399} & 0.010 & 0.036 & 0.371 & 0.570 & 0.752 \\
Breast (human) & \textbf{0.400} & 0.053 & 0.035 & 0.375 & 0.433 & 0.487 \\
Cervical (human) & \textbf{0.178} & 0.017 & 0.014 & 0.157 & 0.247 & 0.026 \\
\textbf{Colon (mouse)} & \textbf{0.615} & -0.137 & 0.072 & 0.588 & 0.739 & 0.744 \\
Colorectal (human) & \textbf{0.431} & 0.056 & 0.083 & 0.419 & 0.532 & 0.623 \\
\textbf{Embryo (mouse)} & \textbf{0.278} & 0.037 & 0.045 & 0.251 & 0.517 & 0.342 \\
Head (zebrafish) & \textbf{0.216} & 0.022 & 0.031 & 0.193 & 0.429 & 0.287 \\
Heart (human) & \textbf{0.292} & 0.010 & 0.021 & 0.266 & 0.688 & 0.304 \\
Kidney (human) & \textbf{0.402} & 0.019 & 0.017 & 0.361 & 0.471 & 0.460 \\
Liver (human) & -0.002 & 0.003 & \textbf{0.010} & -0.002 & 0.561 & 0.400 \\
Lung (human) & \textbf{0.359} & 0.065 & 0.053 & 0.317 & 0.402 & 0.457 \\
Lymph Node (human) & \textbf{0.141} & 0.043 & 0.013 & 0.131 & 0.164 & 0.023 \\
Ovarian (human) & \textbf{0.250} & 0.122 & 0.104 & 0.239 & 0.312 & 0.234 \\
Ovarian glands (human) & \textbf{0.402} & 0.059 & 0.043 & 0.391 & 0.559 & 0.756 \\
Pancreas (human) & \textbf{0.324} & 0.090 & 0.068 & 0.272 & 0.505 & 0.275 \\
Pancreatic (human) & \textbf{0.366} & 0.080 & 0.051 & 0.341 & 0.540 & 0.694 \\
Pancreatic duct gland (human) & \textbf{0.331} & 0.091 & 0.100 & 0.282 & 0.408 & 0.386 \\
Plant (\textit{A. thaliana}) & -0.011 & 0.004 & \textbf{0.012} & -0.009 & 0.385 & 0.052 \\
Prostate (human) & \textbf{0.263} & -0.015 & 0.025 & 0.244 & 0.263 & 0.083 \\
Seed (soybean) & -0.020 & -0.003 & \textbf{0.008} & -0.018 & 0.314 & 0.037 \\
Skin (human) & \textbf{0.359} & 0.048 & 0.086 & 0.299 & 0.437 & 0.368 \\
\textbf{Small Intestine (mouse)} & \textbf{0.572} & -0.104 & 0.067 & 0.548 & 0.701 & 0.715 \\
Tonsil (human) & \textbf{0.294} & 0.028 & 0.018 & 0.249 & 0.462 & 0.399 \\
Xenograft (human + mouse) & \textbf{0.387} & 0.041 & 0.049 & 0.362 & 0.526 & 0.589 \\
\midrule
\textbf{30-sample macro} & \textbf{0.324} & 0.046 & 0.053 &
\textbf{0.291} & \textbf{0.442} & \textbf{0.413} \\
\bottomrule
\end{tabular*}

\endgroup
\end{center}

\textbf{A2 (Independent-sample evidence).}
\begin{itemize}
\item Image is strongest in 21/25 category rows; the four failures remain
visible rather than being removed.
\item Across 18 complete same-organ held-out targets, image
top-50-HVG Gene Pearson averages 0.170 (0.016--0.372), Cell Pearson is 0.275
versus 0.205 for training mean, and F1 is 0.125 versus 0.031. These are sample-,
not donor-held-out, results.
\item The updated benchmark code, fixed splits,
configurations, and evaluation scripts are now available in our
\href{https://anonymous.4open.science/r/sMMC-22M-DB75}{anonymous GitHub
repository}.
\end{itemize}

\textbf{A3 (Visium HD transfer).}\enspace We additionally evaluated eight
complete source-sample-to-target-sample organ pairs. Because these full-panel
transfer results are central to the reviewer's calibration question, we report
them in Q3 rather than only giving a favorable aggregate.

\questionsep
\textbf{Q3.}\quad
\textit{Some gene-wise correlations are low, especially under transfer. How
should they be interpreted biologically, and could marker/cell-type genes or
spatial-only/metadata-only baselines calibrate them?}

\textbf{A1 (Full-panel transfer before calibration).}\enspace The unweighted
per-organ results are:

\begin{center}
\begingroup
\footnotesize
\setlength{\tabcolsep}{3.4pt}
\textit{\textbf{Table 2. Eight-organ full-panel Visium HD transfer before
metadata calibration.} TM denotes the source training mean; its gene-wise
correlation is undefined.}\\[3pt]
\begin{tabular*}{\linewidth}{@{\extracolsep{\fill}}llrrrrr@{}}
\toprule
Organ & Pair & \shortstack{UNI2-h\\Gene P} & \shortstack{UNI2-h\\Cell P} &
\shortstack{UNI2-h\\F1} & \shortstack{TM\\Cell P} & \shortstack{TM\\F1} \\
\midrule
Human breast & j$\to$m & 0.0272 & 0.1109 & 0.1794 & 0.1108 & 0.1215 \\
Human ovary & ad$\to$ak & 0.0267 & 0.4593 & 0.1161 & 0.4906 & 0.0561 \\
Human lung & k$\to$d & 0.0143 & 0.2121 & 0.0960 & 0.1455 & 0.0236 \\
Human pancreas & g$\to$ag & 0.0040 & 0.0090 & 0.0389 & 0.0023 & 0.0168 \\
Human tonsil & u$\to$i & 0.0141 & 0.2769 & 0.0708 & 0.3264 & 0.0346 \\
Mouse brain & e$\to$ah & 0.0243 & 0.2016 & 0.0623 & 0.2238 & 0.0179 \\
Mouse embryo & b$\to$ai & 0.0048 & 0.0913 & 0.0369 & 0.1683 & 0.0134 \\
Mouse kidney & a$\to$aj & 0.0049 & 0.2678 & 0.0513 & 0.4696 & 0.0159 \\
\midrule
\textbf{Organ macro} & --- & \textbf{0.0151} & \textbf{0.2036} &
\textbf{0.0815} & \textbf{0.2422} & \textbf{0.0375} \\
\bottomrule
\end{tabular*}
\endgroup
\end{center}

\begin{itemize}
\item Full-panel Gene Pearson is low, and UNI2-h improves F1 but not Cell
Pearson over the mean. We intentionally retain this negative transfer result.
\item It measures sample, processing, composition, and acquisition shift; it
must not be conflated with the easier training-selected-HVG spatial
interpolation result.
\end{itemize}

\textbf{A2 (Reviewer-inspired transfer calibration).}\enspace We sincerely
thank the reviewer for suggesting metadata/spatial controls. Following this
suggestion, we fit three leakage-controlled baselines on the identical source
cells, target cells, and genes: coordinate-only, segmentation-metadata-only,
and their combination. Segmentation metadata comprise CellViT type/confidence,
polygon/bounding-box morphology, status, and edge flag; expression-derived
\texttt{total\_counts} and \texttt{n\_genes\_detected} are excluded.

\begin{center}
\begingroup
\footnotesize
\setlength{\tabcolsep}{4pt}
\textit{\textbf{Table 3. Segmentation-metadata transfer baseline on the same
eight organ pairs.}}\\[3pt]
\begin{tabular*}{\linewidth}{@{\extracolsep{\fill}}lrrrrr@{}}
\toprule
Organ & Gene P & Gene S & Cell P & Cell S & F1 \\
\midrule
Human breast & 0.0847 & 0.0903 & 0.0785 & 0.1366 & 0.1479 \\
Human ovary & 0.0643 & 0.0638 & 0.4645 & 0.1569 & 0.0615 \\
Human lung & 0.0319 & 0.0314 & 0.1700 & 0.0982 & 0.0267 \\
Human pancreas & 0.0212 & 0.0261 & 0.0043 & 0.0687 & 0.0179 \\
Human tonsil & 0.0282 & 0.0272 & 0.3696 & 0.1432 & 0.0368 \\
Mouse brain & 0.0502 & 0.0467 & 0.2314 & 0.0706 & 0.0223 \\
Mouse embryo & 0.0176 & 0.0193 & 0.1662 & 0.0813 & 0.0146 \\
Mouse kidney & 0.0213 & 0.0211 & 0.3910 & 0.0904 & 0.0175 \\
\midrule
\textbf{Organ macro} & \textbf{0.0399} & \textbf{0.0407} &
\textbf{0.2344} & \textbf{0.1057} & \textbf{0.0432} \\
\bottomrule
\end{tabular*}
\endgroup
\end{center}

\begin{itemize}
\item Segmentation metadata improve Gene Pearson in all 8/8 organs and raise
organ-macro Gene/Cell Pearson from 0.0151/0.2036 to 0.0399/0.2344.
Coordinate-only and metadata+coordinate reach Gene Pearson -0.0023 and 0.0316,
respectively.
\item F1 remains lower than UNI2-h (0.0432 versus 0.0815), so this is a stronger
\emph{correlation-calibration} baseline, not a universally superior predictor.
We will adopt it as the principal non-image transfer baseline alongside the
training mean.
\item Age/sex/disease are not fitted here: each organ has only one source and
one target sample, so clinical-context effects are not separately identifiable.
\end{itemize}

\textbf{A3 (Biological calibration).}\enspace The complementary 30-sample
spatially held-out predictions give:

\begin{center}
\begingroup
\small
\setlength{\tabcolsep}{4pt}
\begin{tabular}{lrrrr}
\toprule
Metric & Image & Coordinate & Spatial KNN & Mean \\
\midrule
Marker/HVG Gene Pearson \(\uparrow\) & 0.201 & 0.031 & 0.028 & N/A \\
Cell-type-stratified pseudobulk RMSE \(\downarrow\) &
0.120 & 0.224 & 0.187 & 0.188 \\
\bottomrule
\end{tabular}
\endgroup
\end{center}

\textbf{A4 (Claim boundary).}
\begin{itemize}
\item The marker inventory overlaps training-selected HVGs, and cell-type strata
are expression-derived; these results show recoverable aggregate structure, not
independent cell-type validation.
\item Under the harder bidirectional all-gene lung/ovary transfer, image Gene
Pearson averages are \(-0.0003/0.0020\), whereas
metadata-only baselines reach 0.104/0.142; coarse-state balanced accuracy is
0.139/0.159 for image prediction versus 0.220/0.252 for metadata.
\item Together, the native-Xenium and eight-organ HD controls demonstrate
context dependence and motivate stronger patient/platform-robust methods
rather than supporting unrestricted cell-wise recovery.
\end{itemize}

\questionsep
\textbf{Q4.}\quad
\textit{In-domain results may reflect local interpolation and spatial
autocorrelation rather than robust generalization.}

\textbf{A.}
\begin{itemize}
\item We will call same-sample results \emph{spatial interpolation diagnostics}.
\item The 30-sample experiment uses contiguous holdouts and a 5\%
train--test buffer, while the 18-target and eight-organ experiments hold out
complete samples. Their lower and heterogeneous performance quantifies the
generalization gap.
\item ``Patient-held-out'' will be used only when donor identity is verified.
\end{itemize}

\questionsep
\textbf{Q5.}\quad
\textit{The ovarian context analysis is narrow and does not establish
systematic evaluation of all variables.}

\textbf{A1 (Scope of the ovarian result).}
\begin{itemize}
\item Ovarian AK (60+) \(\to\) AD (60+) versus AK \(\to\) AL (40-) changes F1
0.289 \(\to\) 0.072, gene Spearman 0.036 \(\to\) 0.006, cell Spearman
0.206 \(\to\) 0.095, and nonzero rate 0.166 \(\to\) 0.027.
\item Because age is nested
within sample/patient and panels differ, we will call this an \emph{age-associated,
sample-confounded shift}, not a causal age effect.
\end{itemize}

\textbf{A2 (Broader context audit).}\enspace To show the broader purpose of rich
context, we applied Average/Worst/Gap audits to complementary cohorts:

\begin{center}
\begingroup
\scriptsize
\setlength{\tabcolsep}{3pt}
\begin{tabularx}{\linewidth}{@{}X p{0.22\linewidth}rrr@{}}
\toprule
Task / encoder & Context; split & Average BA & Worst BA & Gap \\
\midrule
Bone marrow / scGPT & assay; patient-CV & 0.962 & 0.740 & 0.258 \\
Ten tissues / scGPT & dataset; patient-CV & 0.667--0.962 & 0.004--0.892 &
max 0.975 \\
LUAD KRAS / CONCH & site; leave-one-site & 0.499 & 0.375 & 0.542 \\
LGG IDH / H-Optimus-0 & site; leave-one-site & 0.747 & 0.476 & 0.524 \\
\bottomrule
\end{tabularx}
\endgroup
\end{center}

\textbf{A3 (Significance).}\enspace Rich metadata enable subgroup-support and
worst-context audits; they do not themselves remove bias.

\questionsep
\textbf{Q6.}\quad
\textit{Users need clear guidance on which fields are direct measurements and
which are generated or post-processed.}

\textbf{A1 (Provenance classes).}
\begin{center}
\begingroup
\small
\begin{tabularx}{\linewidth}{@{}p{0.14\linewidth}X@{}}
\toprule
Provenance & Examples \\
\midrule
Direct & H\&E; Xenium transcript coordinates/platform masks; Visium HD
2-\(\mu\)m bin counts; source metadata as reported. \\
Processed & Registration; CellViT footprints; Visium HD bin-to-cell matrices;
crops; normalized matrices; splits; molecular-profile summaries. \\
Generated & Metadata-grounded GPT-4o captions and model predictions; neither is
molecular ground truth. \\
\bottomrule
\end{tabularx}
\endgroup
\end{center}

\textbf{A2 (Generated-text validation and scope).}
\begin{itemize}
\item GPT-4o only verbalizes supplied fields.
\item PLIP/CONCH audits give AUC 0.993/0.988 and 100\% Top-3 retrieval.
\item We will restore prompts, field checks, examples, and matrices in the
Appendix.
\item Captions remain optional context, not molecular ground truth.
\end{itemize}

\end{document}
